\chapter{Welfare and the terms of sale}\label{ch:welfare}

\begin{chapterplan}
Who gains when the challenger sees the price, and how do the payment rule and the reserve change what the price
does? Three results of Section~7 of the paper, with proofs in its Appendix~A.7.
\end{chapterplan}

% ====================================================================
\section{Access to prices}\label{sec:welfare-access}

In the \emph{hidden-price economy} the challenger sees neither the price nor the order; market makers still do. With
\(w_j=t_j-t_0\) and class entry \(e_j=\sum_xf_j(x)e(x)\), the terminal payoff \(t_0+e(x)w_j\) of
Chapter~\ref{ch:beliefs} gives, under any entry rule, proceeds \(\mathcal R=t_0+\sum_j\pi_jw_je_j\) and net surplus
\(\mathcal W=v_0+\sum_j\pi_j\{a_je_j-c_H(e_j-\rho)\}\), where \(v_0=\E[\mathcal V_0(R)]\),
\(\mathcal V_0(R)=R\,\1\{R\ge p\}\), and \(a_j=g_j+pF(p)\) is the mean of the next identity.

{\small
\begin{paperresult}{The gain from one prepared challenger (the paper's identity, Appendix A.7)}
For \(\theta\ge p\) and every \(R\ge0\),
\(\max\{\mathcal V_0(R),\theta\}-\mathcal V_0(R)=(\theta-\max\{p,R\})_++p\,\1\{R<p\}\).
\end{paperresult}}
\noindent Both sides equal \(\theta\) if \(R<p\), and \((\theta-R)_+\) if \(R\ge p\).

{\small
\begin{paperresult}{The proposition on access to prices (Appendix A.7 of the paper)}
Hold the incumbent at \(F_1\). Let the hypotheses of Proposition~2 hold. Compare the equilibrium of Proposition~2 with
the equilibrium of an economy in which the challenger does not see the price. All other primitives stay fixed, and
market makers set competitive prices in both economies. Write \(\mathcal R\) for expected target proceeds and
\(\mathcal W\) for expected acquisition surplus net of preparation costs. The superscript \(1\) marks price access
and \(0\) marks the hidden price. Acquisition surplus excludes transfers between bidders, shareholders and traders.
Then both economies have unique full orders. Price access raises expected target proceeds and acquisition surplus net
of preparation costs:
\[
\mathcal R^1-\mathcal R^0=(1-\rho)\bigl\{\pi(b+\lambda)(t_H-t_0)+(1-\pi)b(t_L-t_0)\bigr\}>0,\ 
\mathcal W^1-\mathcal W^0=(1-\rho)(b+\pi\lambda)\bigl[B_1(M)-c_H+pF_1(p)\bigr]>0.
\]
\end{paperresult}}

{\small\begin{derivation}[proof]
Without access the belief is \(\pi\). In every candidate, \(B_1(\pi)\le B_0(\pi)<c_H\) (Proposition~1(a), left side
of (H1)), so \(e\equiv\rho\); and \(K_H\ge\rho\kappa\Delta_1>k\), \(K_L\) alike (right side of (H2)), so orders are
full. With access, entry is \((\rho,\rho,1)\): \(e^1_H-\rho=(1-\rho)(b+\lambda)\), \(e^1_L-\rho=(1-\rho)b\). Insert
in \(\mathcal R\) and \(\mathcal W\); \(t_0\), \(v_0\) and \(c_H\rho\) cancel, and the signs follow from \(w_L\ge0\),
\(w_H=w_L+\Delta_1>0\) and \(B_1(M)>c_H\) (H1). With \(a_j-c_H\) weighted by
\(\pi(b+\lambda)=(b+\pi\lambda)M\) and \((1-\pi)b=(b+\pi\lambda)(1-M)\),
\begin{align*}
\mathcal W^1-\mathcal W^0&=(1-\rho)(b+\pi\lambda)\bigl\{M(a_H-c_H)+(1-M)(a_L-c_H)\bigr\}
=(1-\rho)(b+\pi\lambda)\bigl\{B_1(M)+pF_1(p)-c_H\bigr\}.
\end{align*}
\end{derivation}}

\emph{In words.} A purchase tells a high-cost challenger that its class is probably high, so it prepares; at the
prior it does not. Access adds only these preparations, each worth \(B_1(M)+pF_1(p)-c_H>0\) to the economy.

\begin{watchout}
Each result of this chapter compares equilibria of separate economies, each solved again; none describes one market
over time. Strict comparisons carry the tie-rule claims.
\end{watchout}

% ====================================================================
\section{The payment rule}\label{sec:welfare-payment}

\emph{In words.} With seller weight \(\omega\) and \(\ell<R<h\), class \(H\) pays \((1-\omega)R+\omega h\), and the
incumbent that beats class \(L\) pays \((1-\omega)\ell+\omega R\). The class difference moves with \(R\) at rate
\(1-2\omega\): the second-price rule widens the spread, as Chapter~\ref{ch:ccc} needs.

{\small
\begin{paperresult}{The proposition on the payment rule (Appendix A.7 of the paper)}
Consider a verifiable-value institution with zero reserve. After diligence, values are verifiable, and a sale to the
next buyer at its value is an enforceable fallback. The seller and the winner divide the surplus above the fallback by
Nash bargaining, with seller weight \(\omega\in[0,1)\). Without a challenger, the seller receives \(\omega R\).
Let the challenger's values be the point masses \(\ell\) and \(h\), and let the incumbent distribution have support
\([0,\bar r]\) with \(0<\ell<\bar r<h\). The transfer, the challenger's expected profit and the spread are
\[
T_\omega(R,\theta)=(1-\omega)\min\{R,\theta\}+\omega\max\{R,\theta\},\quad
g_{\theta,\omega}=(1-\omega)\,\E[(\theta-R)_+],\quad
\Delta_\omega=\omega(h-\ell)+(1-2\omega)\,\E[(R-\ell)_+].
\]
A stronger incumbent weakly lowers the challenger's profit. Its effect on the spread weakly has the sign of
\(1-2\omega\). The effect is strict when \(\E[(R-\ell)_+]\) changes strictly and \(\omega\ne1/2\). At \(\omega=1/2\),
the spread equals \((h-\ell)/2\) at every incumbent distribution.
\end{paperresult}}

{\small\begin{derivation}[proof]
The first-order condition of \(\omega\ln(T-a)+(1-\omega)\ln(c-T)\), with \(a=\min\{R,\theta\}\),
\(c=\max\{R,\theta\}\), gives \(T_\omega\). A winner keeps \((1-\omega)(\theta-R)\). For \(R<h\),
\(T_\omega(R,h)-T_\omega(R,\ell)=\omega(h-\ell)+(1-2\omega)(R-\ell)_+\), with mean \(\Delta_\omega\). By the tail
formula, for \(\tilde F\le F\),
\(g_{\theta,\omega}(F)-g_{\theta,\omega}(\tilde F)=(1-\omega)\int_0^\theta(F-\tilde F)\,dy\ge0\) and
\(\Delta_\omega(\tilde F)-\Delta_\omega(F)=(1-2\omega)\int_\ell^{\bar r}(F-\tilde F)\,dy\).
The last integral is nonnegative, and positive exactly when \(\E[(R-\ell)_+]\) changes strictly.
\end{derivation}}

% ====================================================================
\section{Reserves}\label{sec:welfare-reserves}

A reserve with \(\ell<p<h\) excludes the low class. The spread of Chapter~\ref{ch:auction} then has the reserve term
\(pF(p)D(p^-)=pF(p)\), which shrinks with strength. With \(G_H\le G_L\), \(\DT\ge0\), \(B\) is nondecreasing, \(B_1\le B_0\).

{\small
\begin{paperresult}{The proposition on reserves (Appendix A.7 of the paper)}
Fix a reserve, and allow class values below it. Keep the class ranking, the costs, the slot and a strengthening
\(F_1\le F_0\), and compute each class object with the reserve. Suppose \(k<\rho\kappa\min\{\Delta_0,\Delta_1\}\).
\begin{enumerate}[label=(\alph*),leftmargin=2em,nosep]
\item If \(B_0(\pi)<c_H<B_1(M)\), entry is the same at the two strengths:
\(\Ent_0=\Ent_1=\rho+(1-\rho)(b+\pi\lambda)\).
\item If \(B_0(\pi)<c_H<B_0(M)\) and \(B_1(M)<c_H\), entry falls strictly:
\(\Ent_0=\rho+(1-\rho)(b+\pi\lambda)>\rho=\Ent_1\).
\end{enumerate}
Each economy has a unique outcome under every pair of tie rules.
\end{paperresult}}

{\small\begin{derivation}[proof]
By the trading bound, \(K_H,K_L\ge\rho\kappa\Delta_i>k\): orders are full and the three prices differ. Then \(B_0(m)\le B_0(\pi)<c_H\); \(B_0(M)\ge B_1(M)>c_H\) in (a) and \(B_0(M)>c_H\) in (b); and
\(B_1(\mu)\le B_1(M)<c_H\) for all \(\mu\le M\) in (b). So entry is \((\rho,\rho,1)\) and
\(\Ent=\rho+(1-\rho)\Pr(X=1)\), except at \(F_1\) in (b), where \(e\equiv\rho\). All comparisons are strict.
\end{derivation}}

\emph{In words.} With the low class excluded, \(\Delta_0\ge pF_0(p)\) can exceed \(k\), against (H2): the price
already speaks at the weak incumbent, so competition has nothing to switch on.

\begin{chaptersummary}
Price access adds high-cost preparation after a purchase only, and raises proceeds and net surplus. Under Nash
bargaining, a stronger incumbent moves the spread in the direction of \(1-2\omega\). A reserve that excludes the low
class makes entry flat, or makes it fall when \(B_1(M)<c_H\).
\end{chaptersummary}
